
Imagine a scenario in which the members of a group have each ranked a set of alternatives
from most preferred to least preferred.
How should this set of rankings be consolidated to a single ranking
or to a set of winning alternatives that fairly represents the entire group's preferences?
The most familiar setting for this problem is an election,
where voters rank candidates to choose their representatives,
but the problem setup applies much more broadly:
consider
a set of friends deciding what movie to watch or restaurant to visit;
a committee choosing what candidates to shortlist or hire; or
a panel of judges who have ranked competitors in ice skating or ballroom dance.

There are many algorithms (called \textit{voting rules}) for aggregating these group preferences:
how should the group choose which to use?
Deliberation about this decision might involve invoking the details of the algorithms,
or making appeals to axiomatic properties that such algorithms may or may not satisfy.
Some group member might say that Single Transferable Vote (STV) is "more proportional"
than Bloc Plurality for electing committees;
that Instant-Runoff Voting (IRV), unlike Plurality, "avoids the spoiler effect";
or that Plurality exhibits a "center squeeze" effect.
However, non-experts may find such arguments to be overly abstract.
Methods to visualize the behavior of voting rules can bridge this communication gap
and help voters develop intuition about the differences between algorithms.

The core objective of our project is to employ an established election visualization technique,
the two-dimensional spatial model, as the basis of an interactive election simulator.
Allowing users to build, experiment, and iterate upon their own simulations
builds intuition in ways that a set of static plots cannot.
Combining visualization with iterative experimentation
allows users to both \textit{see} the properties of the various algorithms
and to learn the conditions under which these properties emerge.

\textbf{Our Methodology.}
To facilitate visualization, we adopt a two-dimensional spatial model of voters' preferences.
Voters and candidates are jointly embedded in the Euclidean plane,
and voters rank candidates in order of increasing distance.
Voters therefore prefer candidates that are closer, or "more similar", to themselves.
This model is common in both theoretical \cite{spatial}
and empirical \cite{guide} research in computational social science.
Further, it has a natural interpretation:
the position of voters and candidates might represent their
positions on two political issues or on two ideological axes.

In the Spatial Election Explorer (SEE),
voters and candidates can be placed on the plane by hand,
by sampling from the available distributions,
or by mapping survey data to positions on ideological axes.
Their positions determine a unique collection of rankings,
which can be used as input to any of various voting rules.
The winners are displayed, providing the user a visual description
of how the distribution of winners differs according to the choice of voting rule.

\textbf{Audience.}
The intended audience for our application is threefold:
\begin{itemize}[itemsep=0.2em]
\item[1.] Computational social choice researchers can experiment with various voting rules,
testing and refining the intuition they already have for those methods.
Visualization may also suggest patterns of behavior in these voting rules that can
be axiomatized and studied formally.
%
\item[2.] Students studying social choice theory can use the tool to develop intuition
for the behavior of voting rules that supplements their learning in the
classroom.
%
\item[3.] Policymakers and policy advocates considering ranked-choice voting
   can use SEE to compare voting rules and build intuition about how outcomes vary.
\end{itemize}

\textbf{Related Work.}
Our work is primarily inspired by Elkind et al. \cite{elkind},
whose experiments generated plots capturing the distribution of winning candidates
under various multi-winner voting rules
and across varying underlying distributions of voters and candidates.
SEE extends this approach by making it interactive and incorporating real-world survey data.

Another visualization approach is presented by
Szufa et al. \cite{dap}, who created visual "maps" of preference profiles
by computing swap distances between ballots and applying the
multidimensional scaling (MDS) algorithm.
Extending this method to an interactive application would require
some external method for generating election data;
the spatial model, on the other hand, consolidates the generation and visualization steps
because the position of points \textit{determines} the voters' rankings.
